\documentclass[11pt]{article}
\usepackage{manual_docs_style}

\begin{document}

\docTitle{Stage: Final Artifact Scoring}
\phantomsection\label{docstart:final_artifact_scoring_stage}

The responsibility of this stage is to score the final artifacts produced by an agentic rollout.
It reads the persisted terminal-bench run, evaluates each trial's final answer or rule artifact, and writes per-trial and run-level scoring summaries.
We are using the metrics top-1 accuracy and shortlist score, following the ranked-label convention used in the main paper.

\section*{Main Driver}
\begin{DocCode}
python workflows/rollout/evaluate_artifact.py [<agentic_run_id>] [--path PATH]
\end{DocCode}
Note that only one of \code{<agentic_run_id>} or \code{--path} may be passed.
The \code{agentic_run_id} resolves to:
\begin{DocCode}
terminal-bench/runs/<agentic_run_id>
\end{DocCode}

\section*{Inputs}
The stage consumes the rollout output directory and its run-local metadata:
\begin{itemize}
  \item \code{run_metadata.json}, including the dataset path and run identifiers
  \item \code{results.json}, including the terminal-bench trial entries
  \item each trial's \code{scenario_info.json}
  \item each trial's final artifact, either \code{results.json} for direct prediction or \code{rule.py} for code-rule prediction
  \item persisted artifacts such as \code{rule.py}, when the evaluation mode is code-based
\end{itemize}
can you
\section*{Internal Steps}
\begin{enumerate}
  \item Resolve and validate \code{terminal-bench/runs/<agentic_run_id>}.
  \item Load \code{run_metadata.json} and \code{results.json}.
  \item Build a question index from the run-local metadata and aliases in \code{results.json}.
  \item Iterate through each trial entry and load its \code{scenario_info.json}.
  \item Score according to \code{scenario_info.json::eval_mode}:
  \begin{itemize}
    \item \code{direct}: require \code{results.json}, validate that it maps each \code{<filename>.parquet} path to a ranked list of labels, and compare each list with the test label.
    It case \code{results.json} does not follow the convention, then try to allow for flexible parsing. By
    \begin{itemize}
      \item Trying to check whether the format is \code{<filename>.parquet} instead of \code{<filename>.parquet}
    \end{itemize}
    \item \code{code}: require \code{rule.py}, call \code{predict(df) -> list[str]} on each test sample and any held-out \code{other} samples, and compare each returned ranked list with test label.
    Internally \code{evaluate_rule} works as follows:
    \begin{itemize}
      \item For each uuid in \code{sample_manifest.json::<shot_slug>::test_samples} and then after for item in \code{sample_manifest.json::<shot_slug>::others}, create \code{uuid_path}
      \item Call:
      \begin{DocCode}
evaluate_rule(
    uuid_path,
    rule.py,
    model_record_path,
    desc_level=scenario_info.json::desc_level,
    noise_level=scenario_info.json::noise_level,
    evaluation_type="shortlist_score",
)
      \end{DocCode}
      Note that in case of any error the \code{error} field is recorded.
    \end{itemize}
  \end{itemize}
  \item Write each evaluated trial's \code{scores.json}.
\end{enumerate}
If the required artifact is missing or invalid, we are skipping it (not including in the results)

\section*{Materialized Output}
Per-trial scores are stored under each trial directory:
\begin{DocCode}
terminal-bench/runs/<agentic_run_id>/<task_id>/<trial_name>/scores.json
\end{DocCode}

\section*{\code{scores.json} Schema}
Each \code{scores.json} payload is validated by \code{shared/scores_schema.py}.
The required top-level keys are:
\begin{itemize}
  \item \code{agent_run_id}: non-empty run identifier
  \item \code{is_correct_format}: bool, present when \code{scenario_info.json::eval_mode=="direct"}. 
  True if \code{results.json} includes all samples, and if its format matches the \code{<filename>} structure.
  \item \code{final_metric_test}: 
  \begin{itemize}
  \item \code{average_top1_accuracy}: mean top-1 accuracy over test samples, float between 0 and 1.
  \item \code{average_shortlist_score}: mean shortlist score over test samples, float between 0 and 1.
  \item \code{average_num_answers}: mean number of unique labels returned per test sample.
  \end{itemize}
  \item \code{sample_results}: dict with \code{predictions}, \code{top1_correct}, \code{shortlist_score}, \code{num_answers}, \code{sample_type}, \code{error}.
  Here \code{predictions} is the returned ranked-label list, \code{top1_correct} records whether the top-1 prediction is correct, \code{shortlist_score} is the per-sample shortlist score, and \code{num_answers} is the number of unique labels returned.
  \code{sample_type} can be either \code{test} or \code{other}; \code{error} is a string only when scoring that sample raised an error.
  \item \code{final_metric_other}: present when \code{scenario_info.json::eval_mode=="code"}.
  \begin{itemize}
    \item \code{average_top1_accuracy}: mean top-1 accuracy over held-out \code{other} samples, float between 0 and 1.
    \item \code{average_shortlist_score}: mean shortlist score over held-out \code{other} samples, float between 0 and 1.
    \item \code{average_num_answers}: mean number of unique labels returned per held-out sample.
  \end{itemize}
\end{itemize}
For averaging, \code{error} values are considered as top-1 accuracy \code{0}, shortlist score \code{0}, and \code{num_answers=null}, so they are skipped in the \code{average_num_answers} calculation.

\end{document}
